% ============================================================
%  cap_misure.tex — Matemagica
%  Misure: lunghezza, massa, capacità, superficie, volume
%  Equivalenze e misure agrarie
% ============================================================

\chapter{Le misure}

% --- Obiettivi di apprendimento --------------------------------
\section*{Obiettivi di apprendimento}
\begin{itemize}[noitemsep]
    \item Conoscere le unità di misura fondamentali per lunghezza, massa e capacità.
    \item Saper utilizzare multipli e sottomultipli di ciascuna unità.
    \item Comprendere la struttura decimale del Sistema Internazionale di misura.
    \item Applicare la tavola delle misure per convertire tra unità diverse.
    \item Conoscere le misure di superficie e di volume e la loro relazione con le misure lineari.
    \item Comprendere la relazione tra decimetro cubo e litro.
    \item Conoscere le misure agrarie (ara ed ettaro).
\end{itemize}

% ============================================================
\section{Il Sistema Internazionale di misura}
% ============================================================

\begin{definizione}
    Il \textbf{Sistema Internazionale di misura} (abbreviato \textbf{SI}) è il sistema di unità di misura adottato a livello mondiale in ambito scientifico e tecnico. È un sistema \textbf{decimale}: ogni unità è dieci volte più grande di quella immediatamente inferiore.
\end{definizione}

I prefissi del Sistema Internazionale seguono sempre lo stesso schema, indipendentemente dall'unità di misura considerata. Conoscere i prefissi significa saper leggere tutte le misure del sistema.

\begin{center}
    \begin{tabular}{llcc}
        \toprule
        \textbf{Prefisso} & \textbf{Simbolo} & \textbf{Significato}    & \textbf{Valore} \\
        \midrule
        chilo             & k                & mille volte l'unità     & $\times 1000$   \\
        etto              & h                & cento volte l'unità     & $\times 100$    \\
        deca              & da               & dieci volte l'unità     & $\times 10$     \\
        —                 & —                & \textit{unità base}     & $1$             \\
        deci              & d                & un decimo dell'unità    & $\div 10$       \\
        centi             & c                & un centesimo dell'unità & $\div 100$      \\
        milli             & m                & un millesimo dell'unità & $\div 1000$     \\
        \bottomrule
    \end{tabular}
\end{center}

% ============================================================
\section{La tavola delle misure}
% ============================================================

\begin{definizione}
    La \textbf{tavola delle misure} è uno strumento che dispone in ordine decrescente tutti i multipli e sottomultipli di un'unità, da sinistra (valori maggiori) a destra (valori minori). Ogni passaggio verso destra corrisponde a dividere per 10; ogni passaggio verso sinistra corrisponde a moltiplicare per 10.
\end{definizione}

\begin{center}
    \begin{tikzpicture}[font=\small]
        % Celle della tavola
        \foreach \i/\lbl/\sub in {
                0/k/chilo,
                1/h/etto,
                2/da/deca,
                3/{u.b.}/unità,
                4/d/deci,
                5/c/centi,
                6/m/milli}
            {
                \node[draw=blu, fill=azzurro, minimum width=1.6cm, minimum height=0.9cm,
                    inner sep=2pt, align=center] at (\i*1.7,0.5) {\textbf{\lbl}\\[1pt]\tiny\sub};
            }
        % Frecce
        \draw[->, thick, blu] (0.8,-0.3) -- (10.4,-0.3)
        node[midway, below=2pt]{\footnotesize $\div 10$ per ogni passo verso destra};
        \draw[->, thick, blu!60] (10.4,-0.9) -- (0.8,-0.9)
        node[midway, below=2pt]{\footnotesize $\times 10$ per ogni passo verso sinistra};
    \end{tikzpicture}
\end{center}

\begin{regola}
    Per convertire un'unità in un'unità più piccola (verso destra nella tavola), si \textbf{moltiplica} per una potenza di 10. Per convertire in un'unità più grande (verso sinistra), si \textbf{divide} per una potenza di 10. Il numero di zeri da aggiungere o di cifre da spostare corrisponde al numero di passi effettuati nella tavola.
\end{regola}

% ============================================================
\section{Misure di lunghezza}
% ============================================================

\begin{definizione}
    L'\textbf{unità di misura fondamentale della lunghezza} nel SI è il \textbf{metro} (simbolo: \textbf{m}). La lunghezza esprime la distanza tra due punti.
\end{definizione}

\subsection{Multipli e sottomultipli del metro}

\begin{center}
    \begin{tabular}{lllr}
        \toprule
        \textbf{Nome}  & \textbf{Simbolo} & \textbf{Relazione con il metro}   & \textbf{Valore in m} \\
        \midrule
        chilometro     & km               & $1\,\text{km} = 1000\,\text{m}$   & $1000$               \\
        ettometro      & hm               & $1\,\text{hm} = 100\,\text{m}$    & $100$                \\
        decametro      & dam              & $1\,\text{dam} = 10\,\text{m}$    & $10$                 \\
        \textbf{metro} & \textbf{m}       & \textit{unità base}               & $1$                  \\
        decimetro      & dm               & $1\,\text{dm} = {,}1\,\text{m}$   & ${,}1$               \\
        centimetro     & cm               & $1\,\text{cm} = {,}01\,\text{m}$  & ${,}01$              \\
        millimetro     & mm               & $1\,\text{mm} = {,}001\,\text{m}$ & ${,}001$             \\
        \bottomrule
    \end{tabular}
\end{center}

\begin{esempio}
    \textbf{Conversione di lunghezze}

    \medskip
    \textbf{a)} Convertire $3{,}5\,\text{km}$ in metri.

    Da km a m si fanno \textbf{3 passi verso destra} nella tavola, quindi si moltiplica per $10^3 = 1000$:
    \[
        3{,}5\,\text{km} = 3{,}5 \times 1000\,\text{m} = 3500\,\text{m}
    \]

    \medskip
    \textbf{b)} Convertire $47\,\text{cm}$ in metri.

    Da cm a m si fanno \textbf{2 passi verso sinistra}, quindi si divide per $10^2 = 100$:
    \[
        47\,\text{cm} = 47 \div 100\,\text{m} = {,}47\,\text{m}
    \]

    \medskip
    \textbf{c)} Convertire $2\,\text{km}$ e $340\,\text{m}$ in centimetri.

    Prima si converte tutto in metri: $2\,\text{km} + 340\,\text{m} = 2340\,\text{m}$.\\
    Da m a cm si fanno \textbf{2 passi verso destra}, quindi si moltiplica per $100$:
    \[
        2340\,\text{m} = 2340 \times 100\,\text{cm} = 234\,000\,\text{cm}
    \]
\end{esempio}

\begin{attenzione}
    \textbf{Errore comune — confondere dam e dm.}

    Il \textbf{decametro} (dam) è \emph{dieci volte} il metro; il \textbf{decimetro} (dm) è \emph{un decimo} del metro. I simboli si assomigliano, ma il loro valore è molto diverso: $1\,\text{dam} = 100\,\text{dm}$. Quando si legge un simbolo, occorre verificare se la lettera iniziale è \emph{d} (deci, $\div 10$) oppure \emph{da} (deca, $\times 10$).
\end{attenzione}

% ============================================================
\section{Misure di massa}
% ============================================================

\begin{definizione}
    L'\textbf{unità di misura fondamentale della massa} nel SI è il \textbf{grammo} (simbolo: \textbf{g}). La massa indica la quantità di materia contenuta in un corpo.
\end{definizione}

\begin{attenzione}
    \textbf{Massa e peso non sono la stessa cosa.}

    Nel linguaggio comune si usa spesso la parola \emph{peso} per indicare la massa di un oggetto. In fisica, il \textbf{peso} è la forza con cui la Terra attrae un corpo e dipende dalla gravità; la \textbf{massa} è invece una proprietà intrinseca del corpo e non cambia al variare della gravità. In questo capitolo utilizziamo il termine \emph{massa}, come previsto dal SI.
\end{attenzione}

\subsection{Multipli e sottomultipli del grammo}

\begin{center}
    \begin{tabular}{lllr}
        \toprule
        \textbf{Nome}   & \textbf{Simbolo} & \textbf{Relazione con il grammo}  & \textbf{Valore in g} \\
        \midrule
        chilogrammo     & kg               & $1\,\text{kg} = 1000\,\text{g}$   & $1000$               \\
        ettogrammo      & hg               & $1\,\text{hg} = 100\,\text{g}$    & $100$                \\
        decagrammo      & dag              & $1\,\text{dag} = 10\,\text{g}$    & $10$                 \\
        \textbf{grammo} & \textbf{g}       & \textit{unità base}               & $1$                  \\
        decigrammo      & dg               & $1\,\text{dg} = {,}1\,\text{g}$   & ${,}1$               \\
        centigrammo     & cg               & $1\,\text{cg} = {,}01\,\text{g}$  & ${,}01$              \\
        milligrammo     & mg               & $1\,\text{mg} = {,}001\,\text{g}$ & ${,}001$             \\
        \bottomrule
    \end{tabular}
\end{center}

\begin{sapevi}
    \textbf{La tonnellata}

    Nella vita quotidiana si usa spesso la \textbf{tonnellata} (simbolo: t) per indicare masse molto grandi. Non è un'unità del SI, ma è accettata nell'uso corrente: $1\,\text{t} = 1000\,\text{kg} = 1\,000\,000\,\text{g}$.
\end{sapevi}

\begin{esempio}
    \textbf{Conversione di masse}

    \medskip
    \textbf{a)} Convertire $2\,\text{kg}$ e $350\,\text{g}$ in grammi.

    $1\,\text{kg} = 1000\,\text{g}$, quindi:
    \[
        2\,\text{kg} + 350\,\text{g} = 2000\,\text{g} + 350\,\text{g} = 2350\,\text{g}
    \]

    \medskip
    \textbf{b)} Convertire $4500\,\text{mg}$ in grammi.

    Da mg a g si fanno \textbf{3 passi verso sinistra}, quindi si divide per $1000$:
    \[
        4500\,\text{mg} = 4500 \div 1000\,\text{g} = 4{,}5\,\text{g}
    \]
\end{esempio}

% ============================================================
\section{Misure di capacità}
% ============================================================

\begin{definizione}
    L'\textbf{unità di misura fondamentale della capacità} è il \textbf{litro} (simbolo: \textbf{l} oppure \textbf{L}). La capacità indica il volume di liquido o gas che un contenitore può accogliere.
\end{definizione}

\begin{attenzione}
    \textbf{Il simbolo del litro}

    Il litro si scrive con la lettera elle minuscola (\textbf{l}) o, per evitare confusione con il numero 1, con la lettera \textbf{L} maiuscola. Entrambe le forme sono accettate; in questo libro si usa \textbf{l}.
\end{attenzione}

\subsection{Multipli e sottomultipli del litro}

\begin{center}
    \begin{tabular}{lllr}
        \toprule
        \textbf{Nome}  & \textbf{Simbolo} & \textbf{Relazione con il litro}   & \textbf{Valore in l} \\
        \midrule
        chilolitro     & kl               & $1\,\text{kl} = 1000\,\text{l}$   & $1000$               \\
        ettolitro      & hl               & $1\,\text{hl} = 100\,\text{l}$    & $100$                \\
        decalitro      & dal              & $1\,\text{dal} = 10\,\text{l}$    & $10$                 \\
        \textbf{litro} & \textbf{l}       & \textit{unità base}               & $1$                  \\
        decilitro      & dl               & $1\,\text{dl} = {,}1\,\text{l}$   & ${,}1$               \\
        centilitro     & cl               & $1\,\text{cl} = {,}01\,\text{l}$  & ${,}01$              \\
        millilitro     & ml               & $1\,\text{ml} = {,}001\,\text{l}$ & ${,}001$             \\
        \bottomrule
    \end{tabular}
\end{center}

\begin{esempio}
    \textbf{Conversione di capacità}

    \medskip
    \textbf{a)} Convertire $1{,}5\,\text{l}$ in millilitri.

    Da l a ml si fanno \textbf{3 passi verso destra}, quindi si moltiplica per $1000$:
    \[
        1{,}5\,\text{l} = 1{,}5 \times 1000\,\text{ml} = 1500\,\text{ml}
    \]

    \medskip
    \textbf{b)} Convertire $75\,\text{cl}$ in litri.

    Da cl a l si fanno \textbf{2 passi verso sinistra}, quindi si divide per $100$:
    \[
        75\,\text{cl} = 75 \div 100\,\text{l} = {,}75\,\text{l}
    \]
\end{esempio}

% ============================================================
\section{Misure di superficie}
% ============================================================

\begin{definizione}
    La \textbf{superficie} (o \textbf{area}) misura l'estensione di una figura piana. L'unità di misura fondamentale della superficie nel SI è il \textbf{metro quadrato} (simbolo: \textbf{m\textsuperscript{2}}), che corrisponde all'area di un quadrato con il lato di $1\,\text{m}$.
\end{definizione}

\subsection{Perché si moltiplica per 100 (e non per 10)?}

Passando da un'unità lineare alla corrispondente unità di superficie, il fattore di conversione \textbf{si eleva al quadrato}.

\begin{center}
    \begin{tikzpicture}
        % Quadrato 1m x 1m
        \draw[thick, blu, fill=azzurro!60] (0,0) rectangle (3,3);
        \node at (1.5,1.5) {$1\,\text{m}^2$};
        \draw[<->, blu] (0,-0.4) -- (3,-0.4) node[midway,below]{$1\,\text{m}$};
        \draw[<->, blu] (3.4,0) -- (3.4,3) node[midway,right]{$1\,\text{m}$};
        % Divisione in dm
        \foreach \x in {0.3,0.6,...,2.7} \draw[gray!50] (\x,0)--(\x,3);
        \foreach \y in {0.3,0.6,...,2.7} \draw[gray!50] (0,\y)--(3,\y);
        \node[below=1.2cm] at (1.5,0) {\footnotesize $1\,\text{m} = 10\,\text{dm}$, quindi $1\,\text{m}^2 = 10 \times 10\,\text{dm}^2 = 100\,\text{dm}^2$};
    \end{tikzpicture}
\end{center}

\begin{regola}
    Tra unità di superficie consecutive nella tavola, il fattore di conversione è $\mathbf{100}$ (non 10), perché la superficie ha \textbf{due dimensioni}.

    \[
        1\,\text{m}^2 = 100\,\text{dm}^2 \qquad 1\,\text{dm}^2 = 100\,\text{cm}^2 \qquad 1\,\text{cm}^2 = 100\,\text{mm}^2
    \]

    Per passare a un'unità più piccola si moltiplica per $100$; per passare a un'unità più grande si divide per $100$.
\end{regola}

\subsection{Tavola delle misure di superficie}

\begin{center}
    \begin{tabular}{lllr}
        \toprule
        \textbf{Nome}           & \textbf{Simbolo}              & \textbf{Relazione}                         & \textbf{Valore in m\textsuperscript{2}} \\
        \midrule
        chilometro quadrato     & km\textsuperscript{2}         & $1\,\text{km}^2 = 1\,000\,000\,\text{m}^2$ & $1\,000\,000$                           \\
        ettometro quadrato      & hm\textsuperscript{2}         & $1\,\text{hm}^2 = 10\,000\,\text{m}^2$     & $10\,000$                               \\
        decametro quadrato      & dam\textsuperscript{2}        & $1\,\text{dam}^2 = 100\,\text{m}^2$        & $100$                                   \\
        \textbf{metro quadrato} & \textbf{m\textsuperscript{2}} & \textit{unità base}                        & $1$                                     \\
        decimetro quadrato      & dm\textsuperscript{2}         & $1\,\text{dm}^2 = 0{,}01\,\text{m}^2$      & ${,}01$                                 \\
        centimetro quadrato     & cm\textsuperscript{2}         & $1\,\text{cm}^2 = 0{,}0001\,\text{m}^2$    & ${,}0001$                               \\
        millimetro quadrato     & mm\textsuperscript{2}         & $1\,\text{mm}^2 = 0{,}000001\,\text{m}^2$  & ${,}000001$                             \\
        \bottomrule
    \end{tabular}
\end{center}

\begin{esempio}
    \textbf{Conversione di superfici}

    \medskip
    \textbf{a)} Convertire $3\,\text{m}^2$ in centimetri quadrati.

    Da m\textsuperscript{2} a cm\textsuperscript{2} ci sono \textbf{2 passi} nella tavola delle superfici, quindi si moltiplica per $100^2 = 10\,000$:
    \[
        3\,\text{m}^2 = 3 \times 10\,000\,\text{cm}^2 = 30\,000\,\text{cm}^2
    \]

    \medskip
    \textbf{b)} Convertire $450\,\text{dm}^2$ in metri quadrati.

    Da dm\textsuperscript{2} a m\textsuperscript{2} c'è \textbf{1 passo verso sinistra}, quindi si divide per $100$:
    \[
        450\,\text{dm}^2 = 450 \div 100\,\text{m}^2 = 4{,}5\,\text{m}^2
    \]
\end{esempio}

\begin{attenzione}
    \textbf{Errore comune — applicare il fattore lineare alle superfici.}

    Un errore frequente è dividere o moltiplicare per $10$ invece che per $100$ quando si convertono superfici. Ricordare sempre: le superfici hanno due dimensioni, quindi il fattore di conversione è $10^2 = 100$ per ogni passo nella tavola. Ad esempio, $1\,\text{m}^2 \neq 10\,\text{dm}^2$, ma $1\,\text{m}^2 = 100\,\text{dm}^2$.
\end{attenzione}

% ============================================================
\section{Misure di volume}
% ============================================================

\begin{definizione}
    Il \textbf{volume} misura lo spazio occupato da un corpo nello spazio tridimensionale. L'unità di misura fondamentale del volume nel SI è il \textbf{metro cubo} (simbolo: \textbf{m\textsuperscript{3}}), che corrisponde al volume di un cubo con spigolo di $1\,\text{m}$.
\end{definizione}

\subsection{Perché si moltiplica per 1000?}

Passando da un'unità lineare alla corrispondente unità di volume, il fattore di conversione \textbf{si eleva al cubo}.

\begin{center}
    \begin{tikzpicture}
        % Cubo stilizzato in assonometria
        \def\a{2.5}
        \def\s{0.6}
        % Faccia frontale
        \draw[thick, blu, fill=azzurro!60] (0,0) rectangle (\a,\a);
        % Faccia superiore
        \draw[thick, blu, fill=azzurro!30]
        (0,\a) -- (\s,\a+\s) -- (\a+\s,\a+\s) -- (\a,\a) -- cycle;
        % Faccia laterale destra
        \draw[thick, blu, fill=azzurro!45]
        (\a,0) -- (\a+\s,\s) -- (\a+\s,\a+\s) -- (\a,\a) -- cycle;
        % Quote
        \draw[<->, blu] (0,-0.35) -- (\a,-0.35) node[midway,below]{$1\,\text{m}$};
        \draw[<->, blu] (\a+\s+0.1,\s) -- (\a+\s+0.1,\a+\s) node[midway,right]{$1\,\text{m}$};
        \node at (\a/2, \a/2) {$1\,\text{m}^3$};
        \node[below=1.1cm] at (\a/2+0.2,0)
        {\footnotesize $1\,\text{m} = 10\,\text{dm}$, quindi $1\,\text{m}^3 = 10 \times 10 \times 10\,\text{dm}^3 = 1000\,\text{dm}^3$};
    \end{tikzpicture}
\end{center}

\begin{regola}
    Tra unità di volume consecutive nella tavola, il fattore di conversione è $\mathbf{1000}$, perché il volume ha \textbf{tre dimensioni}.

    \[
        1\,\text{m}^3 = 1000\,\text{dm}^3 \qquad 1\,\text{dm}^3 = 1000\,\text{cm}^3 \qquad 1\,\text{cm}^3 = 1000\,\text{mm}^3
    \]

    Per passare a un'unità più piccola si moltiplica per $1000$; per passare a un'unità più grande si divide per $1000$.
\end{regola}

\subsection{Tavola delle misure di volume}

\begin{center}
    \begin{tabular}{lllr}
        \toprule
        \textbf{Nome}       & \textbf{Simbolo}              & \textbf{Relazione}                        & \textbf{Valore in m\textsuperscript{3}} \\
        \midrule
        chilometro cubo     & km\textsuperscript{3}         & $1\,\text{km}^3 = 10^9\,\text{m}^3$       & $1\,000\,000\,000$                      \\
        ettometro cubo      & hm\textsuperscript{3}         & $1\,\text{hm}^3 = 10^6\,\text{m}^3$       & $1\,000\,000$                           \\
        decametro cubo      & dam\textsuperscript{3}        & $1\,\text{dam}^3 = 1000\,\text{m}^3$      & $1\,000$                                \\
        \textbf{metro cubo} & \textbf{m\textsuperscript{3}} & \textit{unità base}                       & $1$                                     \\
        decimetro cubo      & dm\textsuperscript{3}         & $1\,\text{dm}^3 = 0{,}001\,\text{m}^3$    & ${,}001$                                \\
        centimetro cubo     & cm\textsuperscript{3}         & $1\,\text{cm}^3 = 0{,}000001\,\text{m}^3$ & ${,}000001$                             \\
        millimetro cubo     & mm\textsuperscript{3}         & $1\,\text{mm}^3 = 10^{-9}\,\text{m}^3$    & ${,}000000001$                          \\
        \bottomrule
    \end{tabular}
\end{center}

\begin{esempio}
    \textbf{Conversione di volumi}

    \medskip
    \textbf{a)} Convertire $2\,\text{m}^3$ in decimetri cubi.

    Da m\textsuperscript{3} a dm\textsuperscript{3} c'è \textbf{1 passo verso destra}, quindi si moltiplica per $1000$:
    \[
        2\,\text{m}^3 = 2 \times 1000\,\text{dm}^3 = 2000\,\text{dm}^3
    \]

    \medskip
    \textbf{b)} Convertire $5\,\text{dm}^3$ in centimetri cubi.

    Da dm\textsuperscript{3} a cm\textsuperscript{3} c'è \textbf{1 passo verso destra}, quindi si moltiplica per $1000$:
    \[
        5\,\text{dm}^3 = 5 \times 1000\,\text{cm}^3 = 5000\,\text{cm}^3
    \]
\end{esempio}

\begin{attenzione}
    \textbf{Errore comune — applicare il fattore lineare o quadratico ai volumi.}

    Per i volumi il fattore per ogni passo è $10^3 = 1000$, non $10$ né $100$. Questo vale per ogni coppia di unità adiacenti: $1\,\text{dm}^3 = 1000\,\text{cm}^3$, non $100\,\text{cm}^3$.
\end{attenzione}

% ============================================================
\section{Relazione tra decimetro cubo e litro}
% ============================================================

\begin{definizione}
    Il \textbf{decimetro cubo} e il \textbf{litro} sono due unità di misura equivalenti:
    \[
        \boxed{1\,\text{dm}^3 = 1\,\text{l}}
    \]
    In modo analogo: $1\,\text{cm}^3 = 1\,\text{ml}$ e $1\,\text{m}^3 = 1000\,\text{l}$.
\end{definizione}

Questa relazione deriva dalla definizione storica del litro: il litro fu definito come il volume occupato da un chilogrammo di acqua pura a $4\,°\text{C}$, che corrisponde esattamente a un decimetro cubo.

\begin{center}
    \begin{tikzpicture}
        % Cubo dm
        \def\a{1.8}
        \def\s{0.45}
        \draw[thick, blu, fill=azzurro!60] (0,0) rectangle (\a,\a);
        \draw[thick, blu, fill=azzurro!30]
        (0,\a) -- (\s,\a+\s) -- (\a+\s,\a+\s) -- (\a,\a) -- cycle;
        \draw[thick, blu, fill=azzurro!45]
        (\a,0) -- (\a+\s,\s) -- (\a+\s,\a+\s) -- (\a,\a) -- cycle;
        \draw[<->, blu] (0,-0.3) -- (\a,-0.3) node[midway,below]{\footnotesize $1\,\text{dm}$};
        \node at (\a/2,\a/2) {\footnotesize $1\,\text{dm}^3$};
        % Simbolo uguale
        \node at (\a+\s+0.7,\a/2+\s/2) {\Large $=$};
        % Etichetta litro
        \node[draw=blu, fill=azzurro!40, minimum width=1.5cm, minimum height=\a cm+0.2cm,
            inner sep=4pt, align=center] at (\a+\s+1.9,\a/2+\s/2)
        {\footnotesize $1\,\text{l}$};
    \end{tikzpicture}
\end{center}

\begin{regola}
    Grazie alla relazione $1\,\text{dm}^3 = 1\,\text{l}$, è possibile convertire direttamente tra misure di volume e misure di capacità:
    \[
        1\,\text{m}^3 = 1000\,\text{l} \qquad\qquad 1\,\text{cm}^3 = 1\,\text{ml}
    \]
\end{regola}

\begin{esempio}
    \textbf{Conversione tra volume e capacità}

    \medskip
    \textbf{a)} Una vasca da bagno ha un volume di ${,}3\,\text{m}^3$. Quanti litri contiene?

    $1\,\text{m}^3 = 1000\,\text{l}$, quindi:
    \[
        {,}3\,\text{m}^3 = 0{,}3 \times 1000\,\text{l} = 300\,\text{l}
    \]

    \medskip
    \textbf{b)} Una bottiglia contiene $750\,\text{ml}$. Qual è il suo volume in centimetri cubi?

    $1\,\text{ml} = 1\,\text{cm}^3$, quindi:
    \[
        750\,\text{ml} = 750\,\text{cm}^3
    \]
\end{esempio}

\begin{sapevi}
    \textbf{Perché l'acqua è lo standard?}

    Nell'antichità le misure variavano da città a città. Nel 1795, dopo la Rivoluzione Francese, la Francia introdusse il sistema metrico per unificare le misure in tutto il Paese. Il litro fu definito in riferimento all'acqua proprio perché è una sostanza facilmente disponibile ovunque e con proprietà fisiche precise e riproducibili.
\end{sapevi}

% ============================================================
\section{Misure agrarie}
% ============================================================

\begin{definizione}
    Le \textbf{misure agrarie} sono unità di misura della superficie utilizzate tradizionalmente per indicare l'estensione di terreni agricoli. Le principali sono:
    \begin{itemize}[noitemsep]
        \item l'\textbf{ara} (simbolo: \textbf{a}), pari a $1\,\text{dam}^2 = 100\,\text{m}^2$;
        \item l'\textbf{ettaro} (simbolo: \textbf{ha}), pari a $1\,\text{hm}^2 = 10\,000\,\text{m}^2$.
    \end{itemize}
\end{definizione}

\begin{center}
    \begin{tabular}{llr}
        \toprule
        \textbf{Unità}                                                                            & \textbf{Equivalenza}            & \textbf{Valore in m\textsuperscript{2}} \\
        \midrule
        ettaro (ha)                                                                               & $1\,\text{ha} = 100\,\text{a}$  & $10\,000$                               \\
        ara (a)                                                                                   & $1\,\text{a} = 100\,\text{m}^2$ & $100$                                   \\
        \midrule
        \multicolumn{2}{l}{$1\,\text{ha} = 100\,\text{a} = 10\,000\,\text{m}^2 = 1\,\text{hm}^2$} &                                                                           \\
        \bottomrule
    \end{tabular}
\end{center}

\begin{regola}
    \[
        1\,\text{ha} = 100\,\text{a} \qquad 1\,\text{a} = 100\,\text{m}^2 \qquad 1\,\text{ha} = 10\,000\,\text{m}^2
    \]
\end{regola}

\begin{esempio}
    \textbf{Conversione con misure agrarie}

    \medskip
    \textbf{a)} Un campo ha un'estensione di $3{,}5\,\text{ha}$. Quanti metri quadrati sono?

    $1\,\text{ha} = 10\,000\,\text{m}^2$, quindi:
    \[
        3{,}5\,\text{ha} = 3{,}5 \times 10\,000\,\text{m}^2 = 35\,000\,\text{m}^2
    \]

    \medskip
    \textbf{b)} Un terreno misura $2500\,\text{m}^2$. Quante are sono?

    $1\,\text{a} = 100\,\text{m}^2$, quindi:
    \[
        2500\,\text{m}^2 = 2500 \div 100\,\text{a} = 25\,\text{a}
    \]
\end{esempio}

\begin{sapevi}
    \textbf{L'ettaro nella vita quotidiana}

    L'ettaro è tuttora l'unità più usata in agricoltura e nel settore immobiliare per i terreni. In Svizzera, i dati sulla superficie agricola sono sistematicamente espressi in ettari. Un campo da calcio standard ha un'area di circa ${,}7\,\text{ha}$.
\end{sapevi}

% ============================================================
\section{Riepilogo}
% ============================================================

\begin{regola}
    \textbf{Riepilogo delle conversioni fondamentali}

    \medskip
    \textbf{Misure lineari} — fattore $\times 10$ per ogni passo verso destra (unità più piccola):
    \[
        1\,\text{km} = 10\,\text{hm} = 100\,\text{dam} = 1000\,\text{m} = 10\,000\,\text{dm} = 100\,000\,\text{cm} = 1\,000\,000\,\text{mm}
    \]

    \medskip
    \textbf{Misure di massa} — fattore $\times 10$ per ogni passo:
    \[
        1\,\text{kg} = 10\,\text{hg} = 100\,\text{dag} = 1000\,\text{g} = 10\,000\,\text{dg} = 100\,000\,\text{cg} = 1\,000\,000\,\text{mg}
    \]

    \medskip
    \textbf{Misure di capacità} — fattore $\times 10$ per ogni passo:
    \[
        1\,\text{kl} = 10\,\text{hl} = 100\,\text{dal} = 1000\,\text{l} = 10\,000\,\text{dl} = 100\,000\,\text{cl} = 1\,000\,000\,\text{ml}
    \]

    \medskip
    \textbf{Misure di superficie} — fattore $\times 100$ per ogni passo:
    \[
        1\,\text{km}^2 = 100\,\text{hm}^2 = 10\,000\,\text{dam}^2 = \ldots
    \]

    \medskip
    \textbf{Misure di volume} — fattore $\times 1000$ per ogni passo:
    \[
        1\,\text{m}^3 = 1000\,\text{dm}^3 = 1\,000\,000\,\text{cm}^3 = 1\,000\,000\,000\,\text{mm}^3
    \]

    \medskip
    \textbf{Relazione volume–capacità}:
    \[
        1\,\text{dm}^3 = 1\,\text{l} \qquad 1\,\text{cm}^3 = 1\,\text{ml} \qquad 1\,\text{m}^3 = 1000\,\text{l}
    \]

    \medskip
    \textbf{Misure agrarie}:
    \[
        1\,\text{ha} = 100\,\text{a} = 10\,000\,\text{m}^2
    \]
\end{regola}
